\documentclass[10pt,a4paper]{article}
\usepackage{nexusS5}
\setnexusheader{Entrée en Quatrième — Stage de pré-rentrée}{Évaluation finale — Ines KEFI}
\begin{document}
\nexuscover{ÉVALUATION FINALE --- DIAGNOSTIC DE SORTIE}{Durée : 45 minutes}{Ines KEFI}{Entrée en Quatrième}{Mathématiques --- à traiter individuellement}
\par\noindent\begin{tabularx}{\linewidth}{>{\bfseries}p{22mm} X >{\bfseries}p{16mm} X}
\toprule Nom et prénom & \rule{62mm}{0.32pt} & Date & \rule{34mm}{0.32pt} \\ \bottomrule
\end{tabularx}

\begin{goalbox}
\textbf{Ce document est un diagnostic de sortie, pas un classement.} Il mesure ce qui est
désormais installé, ce qui reste fragile, et il sert à écrire le plan des quatre premières
semaines de la rentrée.
\end{goalbox}

\begin{methodbox}
\textbf{Consignes}
\begin{itemize}
\item Durée : \textbf{45 minutes}. Partie A : environ 10 min. Partie B : environ 20 min. Partie C : environ 15 min.
\item Travailler seul, sans document et sans les cartes de rappel de la séance.
\item Écrire la relation ou la propriété utilisée avant le calcul.
\item Une question peut être laissée de côté puis reprise ; il vaut mieux la laisser vide que d'écrire au hasard.
\item Trois lignes « Certitude » seulement : \conf\ , sur la même échelle qu'en début de stage.
\end{itemize}
\end{methodbox}

\newpage
\phasehead{Partie A --- Automatismes et connaissances essentielles}{environ 10 min}{Questions courtes. Écrire le résultat, sans rédaction longue.}
\begin{enumerate}
\needspace{22mm}
\item[\textbf{A1.}] Calculer $(-8) + 3 - (-5)$.
\worklines{2}
\par\vspace{1.2mm}
\needspace{22mm}
\item[\textbf{A2.}] Calculer $\dfrac{5}{8} - \dfrac{1}{4}$.
\worklines{2}
\par\vspace{1.2mm}
\needspace{22mm}
\item[\textbf{A3.}] Réduire $7x + 4 - 3x - 9$.
\worklines{2}
\par\vspace{1.2mm}
\needspace{22mm}
\item[\textbf{A4.}] Dans un triangle rectangle, un angle aigu mesure $37^\circ$. Calculer l'autre angle aigu.
\worklines{2}
\par\vspace{1.2mm}
\needspace{22mm}
\item[\textbf{A5.}] Ranger dans l'ordre croissant : $-5$ ; $2$ ; $-12$ ; $0$ ; $-8$.
\worklines{2}
\par\vspace{1.2mm}
\needspace{22mm}
\item[\textbf{A6.}] Compléter $\dfrac{2}{3}=\dfrac{\ldots}{15}$, puis écrire le facteur utilisé.
\worklines{2}
\par\vspace{1.2mm}
\end{enumerate}
\certbox{Certitude sur l'ensemble de la partie A :}
\needspace{68mm}\par\vspace{3mm}
\phasehead{Partie B --- Application et raisonnement}{environ 20 min}{Écrire la relation, la propriété ou la règle avant le calcul, puis contrôler.}
\begin{enumerate}
\needspace{22mm}
\item[\textbf{B1.}] Une terrasse rectangulaire mesure $12$ m de long et $7$ m de large.
\begin{enumerate}
\item Calculer son aire et son périmètre, en écrivant les unités.
\item Le carrelage coûte $26$ TND le mètre carré. Calculer le prix du carrelage de la terrasse.
\end{enumerate}
\worklines{5}
\par\vspace{1.2mm}
\needspace{22mm}
\item[\textbf{B2.}] \begin{enumerate}
\item Développer $5(x - 3)$.
\item Réduire ensuite $5(x - 3) + 2x + 7$.
\item Contrôler le résultat obtenu pour $x = 2$, en calculant séparément les deux écritures.
\end{enumerate}
\worklines{5}
\par\vspace{1.2mm}
\needspace{22mm}
\item[\textbf{B3.}] Un club compte $40$ membres ; $14$ d'entre eux pratiquent la natation.
\begin{enumerate}
\item Calculer la fréquence des nageurs, en pourcentage.
\item Les âges d'un groupe sont $11$ ; $12$ ; $12$ ; $13$ ; $15$ ; $15$ ; $15$ ; $17$.
Calculer la moyenne, puis vérifier qu'elle est comprise entre la plus petite et la plus grande valeur.
\end{enumerate}
\worklines{5}
\par\vspace{1.2mm}
\needspace{22mm}
\item[\textbf{B4.}] Dans le quadrilatère $ABCD$, les diagonales se coupent en $O$ avec $AO = OC$ et $BO = OD$.
\begin{enumerate}
\item Rédiger en donnée-propriété-conclusion ce que l'on peut affirmer sur $ABCD$.
\item Peut-on affirmer que $ABCD$ est un carré ? Justifier la réponse.
\end{enumerate}
\worklines{6}
\par\vspace{1.2mm}
\end{enumerate}
\certbox{Certitude sur la question B2 :}
\needspace{68mm}\par\vspace{3mm}
\phasehead{Partie C --- Transfert et synthèse}{environ 15 min}{Reconnaître la notion utile, mobiliser plusieurs acquis, conclure par une phrase.}
\begin{enumerate}
\needspace{22mm}
\item[\textbf{C1.}] \textbf{L'atelier de peinture.} Un mur rectangulaire mesure $6$ m de long et $3$ m de haut.
Une fenêtre rectangulaire de $2$ m sur $1{,}5$ m ne doit pas être peinte.
\begin{enumerate}
\item Calculer l'aire de la surface réellement à peindre.
\item Le tiers de cette surface reçoit d'abord une sous-couche. Calculer l'aire concernée.
\item La peinture coûte $19$ TND le pot et le magasin ajoute un forfait de livraison de $25$ TND.
Écrire l'expression $C(n)$ du coût total pour $n$ pots, puis calculer $C(6)$ en montrant la substitution.
\item Le budget est de $200$ TND. Peut-on acheter $10$ pots ? Répondre par un calcul, puis conclure par une phrase.
\end{enumerate}
\worklines{10}
\par\vspace{1.2mm}
\needspace{22mm}
\item[\textbf{C2.}] \begin{enumerate}
\item Un point d'abscisse $-9$ se déplace de $14$ unités vers la droite, puis de $6$ unités vers la gauche.
Donner son abscisse finale.
\item Calculer $(-4)\times(-7)$, puis expliquer le signe obtenu en une phrase.
\end{enumerate}
\worklines{5}
\par\vspace{1.2mm}
\end{enumerate}
\certbox{Certitude sur la question C1 :}
\brouillon{\dimexpr\textheight/3\relax}
\newpage
\sectiontitle{Après l'évaluation --- à remplir en deux minutes}
\begin{methodbox}Ces trois lignes ne sont pas notées. Elles servent à construire le plan des quatre premières semaines de la rentrée.\end{methodbox}
\par\noindent\begin{tabularx}{\linewidth}{>{\bfseries}p{62mm} X}
\toprule
La question que j'ai trouvée la plus sûre & \rule{85mm}{0.32pt} \\
La question sur laquelle j'ai le plus hésité & \rule{85mm}{0.32pt} \\
Une notion que je veux revoir dès la première semaine & \rule{85mm}{0.32pt} \\
\bottomrule\end{tabularx}
\vfill
\begin{graybox}\small Ce document est un diagnostic de sortie. Il sera analysé compétence par compétence, puis comparé au positionnement passé en début de stage lorsque cette comparaison est légitime.\end{graybox}
\end{document}
